%! Composition of Functions — Unit 1, Lesson 1.5 — Lesson Plan
\documentclass[10pt]{article}
\usepackage{precalculus-article}
\usepackage{precalculus-boxes}
\usepackage{graphicx}
\graphicspath{{images/}}
\newcommand{\TallMath}[1]{$\displaystyle #1\rule[-1.4em]{0pt}{3.2em}$}
\newcommand{\CourseName}{Precalculus}
\newcommand{\MeetingLength}{55 minutes}
\newcommand{\UnitNumberName}{Unit 1: Functions and Change \quad}
\newcommand{\LessonNumberName}{Lesson 1.5: Composition of Functions}

\begin{document}
\begin{center}
    {\Large\bfseries \CourseName \\
    {\small\bfseries \UnitNumberName \LessonNumberName}}
\end{center}

\begin{tcolorbox}[colback=lilac, colframe=plum, boxrule=0.9pt, arc=2mm,
    left=2mm, right=2mm, top=1.5mm, bottom=1.5mm]
    \textbf{Primary Objective:} Students will evaluate composite functions
    numerically from a table and analytically from formulas, construct
    $(f \circ g)(x)$ by substituting $g(x)$ for every $x$ in $f(x)$, show that
    $f(g(x))$ and $g(f(x))$ are generally different, and decompose a function
    into an inner and an outer piece --- recognizing a horizontal shift or
    dilation as a composition in disguise.
    \hfill\textit{\small (CED Topic 2.7 --- Skills 1.C, 2.B)}
\end{tcolorbox}

\renewcommand{\arraystretch}{1.6}
\begin{skillbox}[Priority Ideas \& Skills]{goldbox}
\begin{tabularx}{\linewidth}{@{}X X@{}}
\textbf{AP Skills} &
\textbf{Key Understandings} \\
\midrule
\textbf{1.C} --- Construct new functions, using composition, that may be useful
in modeling contexts, criteria, or data. &
\begin{itemize}[leftmargin=1.2em,itemsep=2pt,topsep=0pt]
  \item $(f \circ g)(x) = f(g(x))$: the output of the inner function $g$
        becomes the input of the outer function $f$. (2.7.A.1)
  \item A composite value can be found from a table, a graph, or an analytic
        expression --- always working from the inside out. (2.7.A.2)
\end{itemize} \\
\textbf{2.B} --- Construct equivalent representations of functions and
expressions. &
\begin{itemize}[leftmargin=1.2em,itemsep=2pt,topsep=0pt]
  \item Composition is \textbf{not} commutative: in general
        $f(g(x)) \ne g(f(x))$ --- order changes the answer. (2.7.A.3)
  \item The identity function $f(x) = x$ leaves the other function alone:
        composing with it returns that function unchanged. (2.7.A.4)
\end{itemize} \\
\textbf{Synthesis} --- build, then take apart. &
\begin{itemize}[leftmargin=1.2em,itemsep=2pt,topsep=0pt]
  \item Composition links two quantities that no single formula connects
        directly. (2.7.B.1)
  \item Build $f(g(x))$ analytically by substituting $g(x)$ for \textit{every}
        $x$ in $f(x)$. (2.7.B.2)
  \item Any function can be decomposed into simpler pieces; a horizontal shift
        $f(x+k)$ or dilation $f(kx)$ \textit{is} a composition.
        (2.7.C.1--C.3)
\end{itemize} \\
\end{tabularx}
\end{skillbox}

\begin{skillbox}[{Vocabulary, Concepts, \& Theorems}]{greenbox}
\renewcommand{\arraystretch}{1.4}
\begin{tabularx}{\linewidth}{@{} >{\leavevmode\bfseries\color{plum}}p{0.27\linewidth} X @{}}
Composition            & Feeding one function's output into another function. Written $(f \circ g)(x) = f(g(x))$ and read ``$f$ of $g$ of $x$.'' \\
Inner and outer function & In $f(g(x))$, the \textit{inner} function $g$ runs first; the \textit{outer} function $f$ runs on its result. \\
Identity function      & $I(x) = x$ --- the do-nothing function. $f(I(x)) = f(x)$ and $I(f(x)) = f(x)$. \\
Decomposition          & Writing a single function $h$ as $f(g(x))$ for simpler pieces $f$ and $g$ --- composition run backwards. \\
Domain of a composite  & The inputs $x$ that $g$ accepts \textit{and} whose outputs $g(x)$ are inputs $f$ accepts. If either machine jams, there is no value. \\
\end{tabularx}
\end{skillbox}

\begin{skillbox}[Activate Prior Knowledge \& Spiral Review]{lilac}
    Please see attached warm-up (spiral review).
    Skills reviewed: evaluating a function at a numerical input with the
    parentheses habit (Lesson 1.1), substituting an \textit{expression} rather
    than a number into a rule (Lesson 1.1), reading values off a table both
    directions (Lesson 1.1), and the inside/outside sort from Lesson 1.4.
    Item 1(c) is the load-bearing one --- $f(x^2) = 2x^2+1$ is exactly the
    substitution move Part 3 needs, done once before it has a name.
\end{skillbox}

\begin{skillbox}[Hook]{lilac}
    \textbf{Two registers, one jacket.} A \$60 jacket, a \$10 coupon, and 6\%
    sales tax. Your clerk takes the coupon off first, then adds the tax. The
    clerk at the next register adds the tax first, then takes the coupon off.
    \begin{enumerate}[itemsep=2pt]
        \item Ask for a show of hands: same total, or different? Most classes
              split, and most students argue ``same --- it's the same two
              things.''
        \item Work both totals on the board: \$53.00 and \$53.60. Sixty cents
              apart, from the same two operations.
    \end{enumerate}
    \textit{Discussion prompt:} the coupon and the tax are each a
    \textit{function} of the price. Today we chain functions together, and the
    order of the chain is part of the answer.
\end{skillbox}

\begin{skillbox}[Lesson]{lilac}
\begin{multicols}{2}

\textbf{Part 1: Two Machines in a Row} (8 min)

\begin{itemize}
  \item Draw the two-machine picture: $x$ goes into $g$, $g(x)$ comes out and
        goes straight into $f$, and $f(g(x))$ comes out the far end.
  \item Give the notation once and insist on it: $(f \circ g)(x) = f(g(x))$,
        read ``$f$ of $g$ of $x$.'' The circle is not multiplication.
  \item The organizing sentence, repeated all lesson: \textbf{work from the
        inside out --- the inner function runs first.} In $f(g(x))$ the letter
        written \textit{last} runs \textit{first}.
\end{itemize}

\textbf{Part 2: Composite Values from a Table} (12 min)

\begin{itemize}
  \item Using the notes table, do $f(g(1))$ out loud in two recorded steps:
        first $g(1) = 3$, then $f(3) = 2$. Insist on writing the middle number
        down --- that habit is the whole part.
  \item Then $g(f(1))$: $f(1) = 4$, so $g(4) = 4$. Different answer, same two
        functions. Name it: composition is not commutative (2.7.A.3).
  \item Domain at exposure depth: $f(g(2))$ needs $f(7)$, and 7 is not an input
        in the table --- no value. The inner output has to be something the
        outer function accepts.
\end{itemize}

\columnbreak

\textbf{Part 3: Building a Composite Analytically} (12 min)

\begin{itemize}
  \item State the rule as a recipe: to build $f(g(x))$, \textbf{substitute
        $g(x)$ for every $x$ in $f(x)$} (2.7.B.2). Say ``every.''
  \item Work $f(g(x))$ and then $g(f(x))$ for $f(x) = 2x+1$, $g(x) = x^2-3$.
        The results are $2x^2-5$ and $4x^2+4x-2$ --- not even the same shape.
  \item Check a value both ways: $f(g(2)) = 3$ from the formula and from the
        two-step evaluation. Same number, two routes (2.7.A.2).
\end{itemize}

\textbf{Part 4: Order Matters --- the Checkout Counter} (10 min)

\begin{itemize}
  \item Name the hook's two functions: coupon $c(p) = p - 10$, tax
        $t(x) = 1.06x$. Neither one alone answers ``what do I pay?'' ---
        chaining them does (2.7.B.1).
  \item Build both orders symbolically: $t(c(p)) = 1.06p - 10.60$ and
        $c(t(p)) = 1.06p - 10$. They differ by \$0.60 at \textit{every} price,
        because the second one taxes the ten dollars you never paid.
\end{itemize}

\textbf{Part 5: Decomposition, Identity, and Lesson 1.4} (8 min)

\begin{itemize}
  \item Decompose $h(x) = (3x+2)^4$: inner $g(x) = 3x+2$, outer $f(u) = u^4$.
        Test it by substituting back. ``What would you do first with a
        calculator?'' finds the inner function every time.
  \item Identity function $I(x) = x$: composing with it changes nothing
        (2.7.A.4). Flag it --- it is the target of Lesson 1.6.
  \item \textbf{Deliver on 1.4's promise.} Every $f(x+3)$ they wrote last
        lesson was $f(g(x))$ with $g(x) = x+3$; every $f(2x)$ was $f(g(x))$
        with $g(x) = 2x$. Inside changes were compositions all along ---
        that is why they act on inputs (2.7.C.2--C.3).
\end{itemize}
\end{multicols}
\end{skillbox}

\begin{skillbox}[Active Monitoring]{redbox}
While students work on guided notes and practice, circulate and check:
\begin{itemize}
  \item \textbf{Common error:} evaluating $f(g(1))$ as $f(1)$ times $g(1)$, or
        as $f(1)$ then looking up $g$. Cold-call: ``Which number goes into the
        second machine? Write it down before you use it.''
  \item \textbf{Common error:} running the outer function first --- computing
        $f(1)$ when the problem says $f(g(1))$. Point at the parentheses:
        ``Innermost parentheses first, like any arithmetic.''
  \item \textbf{Common error:} in $f(g(x))$ with $f(x) = 2x+1$, substituting
        into only the first $x$. Ask them to circle every $x$ in $f$ before
        replacing anything.
  \item \textbf{Common error:} asserting $f(g(x)) = g(f(x))$ because ``it's
        like multiplying.'' Send them back to the checkout counter --- \$53.00
        versus \$53.60 settles it faster than any algebra.
  \item \textbf{Common error:} decomposing backwards, calling $f(u) = 3u+2$ the
        outer function of $(3x+2)^4$. Have them substitute back and see they
        built $3x^4+2$ instead.
\end{itemize}
\end{skillbox}

\begin{skillbox}[Group Work \& Differentiation]{redbox}
Students work on the \textbf{Group Activity: The Checkout Counter} in leveled
tiers. Every tier uses the same two functions --- coupon $c(p) = p-10$ and tax
$t(x) = 1.06x$ --- so groups can be regrouped mid-activity without new setup.

\begin{itemize}
  \item \textbf{Tier R (Remediate):} Fill a table of $c(p)$, then a second row
        of $t(c(p))$, one machine at a time; then answer which function ran
        first and what the middle row physically \textit{is}.
  \item \textbf{Tier A (Apply):} Compute $t(c(60))$ and $c(t(60))$ in recorded
        steps, build both formulas symbolically, and state which order the
        shopper wants and by how much.
  \item \textbf{Tier E (Extend):} Replace the coupon with a 25\% discount
        $d(p) = 0.75p$ and discover that $d$ and $t$ \textit{do} commute ---
        then explain why. Finish by decomposing the register's printed rule and
        rewriting $f(x+3)$ from Lesson 1.4 as a composition.
\end{itemize}
\end{skillbox}

\begin{skillbox}[Individual Work \& Assessment]{redbox}
\textbf{Exit Ticket} (last 5--7 minutes, individual, collected):
\begin{enumerate}
  \item For $f(x) = x+4$ and $g(x) = 3x$, evaluate $f(g(2))$ and $g(f(2))$
        (work space provided).
  \item For the same $f$ and $g$, write $(f \circ g)(x)$ and $(g \circ f)(x)$
        simplified, then say whether they are the same function.
  \item Decompose $h(x) = (x-5)^3$ into an outer function and an inner
        function.
\end{enumerate}
\textbf{Assessment note:} Item 1 is the diagnostic --- two matching answers
means the student ran the same machine twice and Part 2 needs a reteach in
tomorrow's warm-up window. Item 2 misses split two ways: $3x+12$ in both slots
is a substitution error, while $3(x+4)$ left unsimplified is only a bookkeeping
miss. Item 3 with $f(u) = u - 5$ on the outside is the decomposition reversal
from Active Monitoring.
\end{skillbox}

\begin{skillbox}[Reinforcement \& Extension]{goldbox}
\textbf{Homework:} 5 problems covering table-based composition in both orders,
analytic construction of $f \circ g$ and $g \circ f$ for a quadratic and a
linear function, the checkout-counter context at a new price, decomposition
(including rewriting a Lesson 1.4 horizontal shift as a composition), and one
AP-style multiple-choice item on $f \circ g$.

\textbf{Extension (optional):} Verify that the identity $I(x) = x$ leaves
$f(x) = 5x-2$ unchanged in either order, then hunt for a function $g$ with
$f(g(x)) = x$. Students who find $g(x) = \dfrac{x+2}{5}$ have invented the
inverse function a day early.

\textbf{Preview:} Next lesson (1.6, Inverse Functions) answers the question the
extension raises: which function \textit{undoes} $f$? The test is a
composition --- $f(f^{-1}(x)) = x$, the identity function from Part 5.
\end{skillbox}

\begin{teachernote}[Warm-Up]
Five minutes, silent start, no calculators. Item 1(c) is the payoff: substituting
$x^2$ into $f(x) = 2x+1$ is a routine algebra question here and the entire
mechanic of Part 3 --- when you get there, say out loud that they already did
it. Item 2's backwards table read (``which $x$ gives 0?'') is the habit Part 2
leans on when the inner output has to be looked up as an input. Item 3 is a
one-line callback to 1.4 so Part 5's punchline lands.
\end{teachernote}

\begin{teachernote}[Guided Notes]
The single most valuable thing you can enforce today is \textbf{writing the
middle number down}. A student who computes $f(g(1))$ in their head and writes
only the final answer has no place to be wrong out loud, and no place for you to
catch it. Part 2's two-step layout is built for that; do not let them skip the
first blank. Part 4 is the meaning of the lesson, not a detour --- if you are
short on time, compress Part 3's second work block ($g(f(x))$) rather than
cutting Part 4 or 5. Domain of a composite stays at exposure depth: the table
value $f(g(2))$ that has no answer is enough, and no student needs interval
algebra today.
\end{teachernote}

\begin{teachernote}[Group Activity]
All three tiers use the same coupon and tax functions, so a student who finishes
Tier R can move straight into Tier A without re-reading a scenario. Tier R's
middle row is the point of that tier: make groups say what the number
\textit{is} (the price after the coupon, before tax) rather than just computing
it. Tier A's \$0.60 gap is worth a whole-class pause --- ask why it is exactly
sixty cents at every price, and let someone get to ``it's the tax on the ten
dollars.'' Tier E's discovery that a percentage discount \textit{does} commute
with tax is the best share of the day: order matters in general, not always,
and knowing the difference is the sophisticated version of the idea.
\end{teachernote}

\begin{teachernote}[Exit Ticket]
Collected, completion credit only. Sort by item 1: $f(g(2)) = 10$ and
$g(f(2)) = 18$ is the target, and any pile where the two answers match is your
reteach group. On item 2 accept $3(x+4)$ as well as $3x+12$ --- the second slot
is testing substitution, not distribution. Item 3 has more than one defensible
answer if a student splits it differently; the test is whether substituting back
rebuilds $h$, so read their check, not their letters.
\end{teachernote}

\begin{teachernote}[Homework]
Problem 4(b) --- rewriting $f(x+6)$ as a composition --- is the one that closes
the loop with Lesson 1.4, and it is the item to scan first; a student who cannot
do it did not hear Part 5. On problem 5, distractor B is $(g \circ f)(x)$ left
unexpanded, C is $(g \circ f)(x)$ expanded, and D is the ``circle means
multiply'' error, $f(x) \cdot g(x)$ --- each is worth thirty seconds of autopsy
tomorrow. The extension is the direct on-ramp to 1.6; if the class never
reached Tier E, assign it explicitly so someone arrives tomorrow having already
built an inverse.
\end{teachernote}
\end{document}
